\label{ch:zustand_prozess}
% Source document: 247

\epigraph{The mere storage of information is energy-neutral.
Only processing has an energy dimension.}{Doc.\,247}

\section{The Question}

In information physics --- represented by Wheeler (\enquote{It from Bit}),
M.\,M.~Vopson \textbf{[V1--V3]} and others --- information is often treated as an ontological
primitive. Landauer's principle is cited as evidence: Since
erasing a bit costs energy, information is physical.

This conclusion confuses two structurally different concepts.

\begin{keyresult}
\textbf{Information as state} --- a configuration, a stored
bit, a winding number --- exists without energy expenditure. It is a
topological or geometric property.\\[4pt]
\textbf{Information as process} --- reading, writing, erasing, transforming
--- costs energy. Landauer's principle applies \emph{only} to erasure,
i.e., to a process, not to the state itself.
\end{keyresult}

\section{The Permanent Memory as a Counter-Experiment}

Anyone who flattens the difference must accept a direct empirical consequence:
Every stored bit on a switched-off device would
continuously consume energy --- which would make permanent memory fundamentally
impossible.

Flash memory holds millions of bits for years at nearly zero energy consumption.
It is a counter-experiment to any theory that equates information and energy.
The existence of such memories directly refutes the equation,
without any room for interpretation.

\section{FFGFT: States are Topological Configurations}

In FFGFT, states are geometric configurations on $T^4$:
winding numbers $(n_\theta, n_\varphi) \in \mathbb{Z}^2$.

The rest mass of a particle in FFGFT is:
\[
  m = r_i \cdot \xi^{p_i} \cdot v, \qquad v = 246\;\text{GeV}
\]
The winding numbers $(n_\theta, n_\varphi)$ are a state --- they cost
no energy as long as the topology remains intact. Only the transition
$n \to n'$ --- a process --- requires or generates energy.

Static energy ($E_0 = mc^2$) corresponds to topologically stable
winding numbers. It does not contribute to thermodynamic entropy
as long as the topology remains intact. Dynamic kinetic energy
corresponds to motion through $T^4$ --- it is thermodynamic and
entropy-producing.

\section{M.\,M.~Vopson and the FFGFT Restriction}

M.\,M.~Vopson postulates in \textbf{[V1]} a mass-information equivalence: Every bit
carries a universal rest mass, from which a gravitational effect follows.

FFGFT restriction: There is no universal bit mass.
$\Ebit = \hbar c / L$ depends on the physical scale $L$ (Doc.\,257).
At the Planck scale, $\Ebit = E_P$; at the nanometer scale it is
a few meV. Bits on different scales have different energies ---
Vopson's \textbf{[V1]} system-independent bit mass is incompatible with the FFGFT geometry.

\begin{ffgftbox}
\textbf{Wheeler:} \enquote{It from Bit} --- geometry from information.\\
\textbf{M.\,M.~Vopson} {[V1,V2]}: Bit mass universal, information has rest mass.\\
\textbf{FFGFT:} \enquote{Bit from It} --- information as projection from
geometry; bit mass scale-sensitive; storing is free, erasing costs.
\end{ffgftbox}

\section{What Landauer Really Shows}

Landauer's principle is not evidence for the physical reality of
information as an ontological primitive. It is evidence for the physical
reality of the \emph{carrier}: When a carrier is irreversibly reconfigured,
Landauer costs arise. The information on the carrier plays no
role in this --- thermodynamics counts regions, not meanings (Doc.\,A271).